% =====================================================================
% CAP. 11 - NIVEL 3 - QUESTOES GRAFICAS E INTEGRADAS
% =====================================================================
\blocoaprofundamento

\begin{questao}[3]{}
Um capacitor de $\SI{40}{\micro\farad}$ permaneceu ligado por muito tempo ao
circuito (a). Em $t=0$, a chave transfere instantaneamente o capacitor para o
circuito (b). Determine $v_C(t)$ para $t>0$, o instante em que a tensão cruza
zero e a corrente no resistor $R_2$ imediatamente após a comutação, tomando como
positivo o sentido do capacitor para a fonte de $-\SI{5}{\volt}$.

\circuitoqq{%
\begin{circuitikz}[scale=0.88,transform shape,line width=0.8pt,every node/.append style={font=\scriptsize}]
\draw (0,0) to[american voltage source,l_={$\SI{15}{\volt}$}] (0,3)
      to[R,l={$R_1=\SI{5}{\kilo\ohm}$}] (3.2,3)
      to[C,v^>={$v_C$}] (3.2,0) -- (0,0);
\node[anchor=east] at (2.95,1.5) {$C=\SI{40}{\micro\farad}$};
\node[font=\sffamily\bfseries\scriptsize] at (1.6,-0.6) {(a) $t<0$};
\begin{scope}[xshift=7.6cm]
\draw (0,0) to[american voltage source,l_={$-\SI{5}{\volt}$}] (0,3)
      to[R,l={$R_2=\SI{10}{\kilo\ohm}$}] (3.2,3)
      to[C,v^>={$v_C$}] (3.2,0) -- (0,0);
\node[anchor=east] at (2.95,1.5) {$C=\SI{40}{\micro\farad}$};
\node[font=\sffamily\bfseries\scriptsize] at (1.6,-0.6) {(b) $t>0$};
\end{scope}
\end{circuitikz}}{Circuitos equivalentes antes e depois da comutação.}

\resposta{$v_C(t)=-5+20e^{-t/0{,}4}\ \si{\volt}$;
$t_{v=0}\approx\SI{0.555}{\second}$; $i_{R_2}(0^+)=\SI{2}{\milli\ampere}$}
\resolucao{Em regime permanente no circuito (a), o capacitor está aberto e não há
queda em $R_1$, de modo que $v_C(0^-)=\SI{15}{\volt}$. Logo
$v_C(0^+)=\SI{15}{\volt}$. Para $t>0$, o valor final é $-\SI{5}{\volt}$ e
\[
\tau=R_2C=(10\,\mathrm{k\Omega})(40\,\mathrm{\mu F})=\SI{0.4}{\second}.
\]
Portanto
\[
v_C(t)=-5+(15+5)e^{-t/0{,}4}=-5+20e^{-t/0{,}4}.
\]
No cruzamento por zero, $e^{-t/0{,}4}=1/4$, logo
$t=0{,}4\ln4\approx\SI{0.555}{\second}$. A corrente inicial em $R_2$ é
$(15-(-5))/10\,\mathrm{k\Omega}=\SI{2}{\milli\ampere}$.}
\end{questao}

\begin{questao}[3]{}
A chave permaneceu fechada por muito tempo e abre em $t=0$. O indutor passa a
liberar sua energia pela associação $R_2\parallel R_3$. Determine a corrente
$i_L(t)$ para $t>0$ e a energia ainda armazenada no indutor em
$t=\SI{50}{\milli\second}$.

\circuitoqq{%
\begin{circuitikz}[scale=0.92,transform shape,line width=0.8pt,every node/.append style={font=\scriptsize}]
\draw (0,0) to[american voltage source,l_={$\SI{30}{\volt}$}] (0,3)
      to[switch] (1.8,3)
      to[R,l={$R_1=\SI{5}{\ohm}$}] (4.0,3) coordinate(a);
\node[anchor=south] at (0.9,3.30) {abre em $t=0$};
\draw (a) to[L,i>^={$i_L$}] (4.0,0) -- (0,0);
\node[anchor=east] at (3.75,1.5) {$L=\SI{200}{\milli\henry}$};
\draw (a) -- (6.3,3) to[R,l_={$R_2=\SI{8}{\ohm}$}] (6.3,0) -- (4.0,0);
\draw (a) -- (8.6,3) to[R,l_={$R_3=\SI{24}{\ohm}$}] (8.6,0) -- (6.3,0);
\end{circuitikz}}{Descarga de um indutor por uma rede resistiva equivalente.}

\resposta{$i_L(t)=6e^{-t/0{,}0333}\ \si{\ampere}$;
$W_L(50\,\mathrm{ms})\approx\SI{0.179}{\joule}$}
\resolucao{Antes da abertura, o indutor ideal está em curto e os resistores
$R_2$ e $R_3$ ficam sem tensão. Assim $i_L(0^-)=30/5=\SI{6}{\ampere}$ e
$i_L(0^+)=\SI{6}{\ampere}$. Para $t>0$,
\[
R_{\mathrm{eq}}=8\parallel24=\SI{6}{\ohm},\qquad
\tau=\frac{0{,}2}{6}\approx\SI{33.3}{\milli\second}.
\]
Logo $i_L(t)=6e^{-t/\tau}$. Em $\SI{50}{\milli\second}$,
$i_L=6e^{-1{,}5}\approx\SI{1.339}{\ampere}$. A energia remanescente é
\[
W_L=\frac12Li_L^2\approx\frac12(0{,}2)(1{,}339)^2
\approx\SI{0.179}{\joule}.
\]}
\end{questao}

\begin{questao}[3]{}
O capacitor possui $v_C(0^-)=\SI{10}{\volt}$. Em $t=0$ a chave fecha e o conecta
à rede mostrada. Determine $v_C(t)$, a resistência equivalente vista pelo
capacitor e o tempo necessário para que o erro em relação ao valor final seja
menor que \SI{1}{\percent} da variação inicial.

\circuitoqq{%
\begin{circuitikz}[scale=0.91,transform shape,line width=0.8pt,every node/.append style={font=\scriptsize}]
\draw (0,0) to[american voltage source,l_={$\SI{18}{\volt}$}] (0,3)
      to[R,l={$R_1=\SI{6}{\kilo\ohm}$}] (3.2,3) coordinate(b);
\draw (b) to[R,l_={$R_2=\SI{3}{\kilo\ohm}$}] (3.2,0) -- (0,0);
\draw (b) to[R,l={$R_3=\SI{2}{\kilo\ohm}$}] (5.7,3)
      to[switch] (7.2,3)
      to[C,v^>={$v_C$}] (7.2,0) -- (3.2,0);
\node[anchor=south] at (6.45,3.30) {fecha em $t=0$};
\node[anchor=east] at (6.95,1.5) {$C=\SI{25}{\micro\farad}$};
\end{circuitikz}}{Capacitor conectado a uma rede de Thévenin não trivial.}

\resposta{$R_{\mathrm{eq}}=\SI{4}{\kilo\ohm}$; $\tau=\SI{0.10}{\second}$;
$v_C(t)=6+4e^{-10t}\ \si{\volt}$; $t_{1\%}\approx\SI{0.461}{\second}$}
\resolucao{Em regime final não circula corrente em $R_3$, portanto o capacitor
assume a tensão do divisor formado por $R_1$ e $R_2$:
$V_f=18\cdot3/(6+3)=\SI{6}{\volt}$. Para obter a resistência vista pelo capacitor,
a fonte ideal de tensão é curto-circuitada:
\[
R_{\mathrm{eq}}=R_3+(R_1\parallel R_2)=2\,\mathrm{k\Omega}+2\,\mathrm{k\Omega}
=\SI{4}{\kilo\ohm}.
\]
Logo $\tau=R_{\mathrm{eq}}C=\SI{0.10}{\second}$ e
$v_C(t)=6+(10-6)e^{-t/0{,}1}=6+4e^{-10t}$. O erro normalizado é
$e^{-t/\tau}$. Para ficar abaixo de \SI{1}{\percent},
$t>\tau\ln100\approx4{,}605\tau=\SI{0.461}{\second}$.}
\end{questao}

\begin{questao}[3]{}
Para $t<0$, a fonte de corrente de $\SI{2}{\ampere}$ permaneceu conectada por
muito tempo ao RLC em paralelo. Em $t=0$ a chave abre e a fonte é removida. Determine
o tipo de resposta natural e obtenha $v(t)$ para $t>0$, considerando a polaridade
indicada.

\circuitoqq{%
\begin{circuitikz}[scale=0.94,transform shape,line width=0.8pt,every node/.append style={font=\scriptsize}]
\draw (0,0) to[american current source] (0,3)
      to[switch] (1.9,3) -- (7.3,3);
\node[anchor=east] at (-0.25,1.5) {$\SI{2}{\ampere}$};
\node[anchor=south] at (0.95,3.30) {abre em $t=0$};
\draw (2.6,3) to[R] (2.6,0);
\node[anchor=east] at (2.35,1.5) {$R=\SI{10}{\ohm}$};
\draw (4.9,3) to[C,v^>={$v$}] (4.9,0);
\node[anchor=east] at (4.65,1.5) {$C=\SI{200}{\micro\farad}$};
\draw (7.3,3) to[L,i>^={$i_L$}] (7.3,0) -- (0,0);
\node[anchor=east] at (7.05,1.5) {$L=\SI{50}{\milli\henry}$};
\end{circuitikz}}{Resposta natural de um circuito RLC em paralelo com energia inicialmente no indutor.}

\resposta{Subamortecida;
$v(t)\approx-51{,}64e^{-250t}\sin(193{,}65t)\ \si{\volt}$}
\resolucao{Antes da abertura, em regime CC, o capacitor é aberto e o indutor é
curto. Assim $v(0^-)=0$ e $i_L(0^-)=\SI{2}{\ampere}$; por continuidade,
$v(0^+)=0$ e $i_L(0^+)=\SI{2}{\ampere}$. Para o RLC em paralelo,
\[
\alpha=\frac{1}{2RC}=\SI{250}{\per\second},\qquad
\omega_0=\frac{1}{\sqrt{LC}}\approx\SI{316.23}{\radian\per\second},
\]
logo a resposta é subamortecida e
$\omega_d\approx\SI{193.65}{\radian\per\second}$. Pela KCL em $0^+$,
$C\dot v(0^+)+i_L(0^+)=0$, portanto
$\dot v(0^+)=-2/(200\,\mathrm{\mu F})=-\SI{10000}{\volt\per\second}$.
Como $v(0)=0$, a parcela cossenoidal é nula e
$B=\dot v(0)/\omega_d\approx-51{,}64$, resultando na expressão indicada.}
\end{questao}

\begin{questao}[3]{}
O circuito RLC em série parte sem energia armazenada e a chave fecha em $t=0$.
Determine a resposta completa $v_C(t)$ e o valor do primeiro pico da tensão do
capacitor.

\circuitoqq{%
\begin{circuitikz}[scale=0.90,transform shape,line width=0.8pt,every node/.append style={font=\scriptsize}]
\draw (0,0) to[american voltage source,l_={$\SI{10}{\volt}$}] (0,3)
      to[switch] (1.8,3)
      to[R,l={$R=\SI{20}{\ohm}$}] (4.0,3)
      to[L,l={$L=\SI{100}{\milli\henry}$}] (6.5,3)
      to[C,v^>={$v_C$}] (9.0,3)
      -- (9.0,0) -- (0,0);
\node[anchor=south] at (0.9,3.30) {fecha em $t=0$};
\node[anchor=south] at (7.75,3.30) {$C=\SI{100}{\micro\farad}$};
\end{circuitikz}}{Resposta ao degrau de um circuito RLC em série subamortecido.}

\resposta{$v_C(t)=10\left[1-e^{-100t}\left(\cos300t+\frac13\sin300t\right)\right]\,\si{\volt}$;
$v_{C,\mathrm{pico}}\approx\SI{13.51}{\volt}$ em $t_p\approx\SI{10.47}{\milli\second}$}
\resolucao{Temos
$\alpha=R/(2L)=\SI{100}{\per\second}$,
$\omega_0\approx\SI{316.23}{\radian\per\second}$ e
$\omega_d=\SI{300}{\radian\per\second}$. Para entrada degrau e condições iniciais
nulas, a resposta do capacitor é
\[
v_C(t)=V_s\left[1-e^{-\alpha t}\left(\cos\omega_dt+
\frac{\alpha}{\omega_d}\sin\omega_dt\right)\right].
\]
Substituindo os valores, obtém-se a expressão indicada. O primeiro pico ocorre em
$t_p=\pi/\omega_d\approx\SI{10.47}{\milli\second}$. Nesse instante,
$\sin(\omega_dt_p)=0$ e $\cos(\omega_dt_p)=-1$, logo
$v_C(t_p)=10[1+e^{-100t_p}]\approx\SI{13.51}{\volt}$.}
\end{questao}

\begin{questao}[3]{}
Em laboratório, mediu-se a corrente natural do RLC em série mostrado. A distância
entre picos consecutivos é aproximadamente $T_d=\SI{8}{\milli\second}$ e a
envoltória cai à metade a cada $\SI{16}{\milli\second}$. Sabendo que
$R=\SI{8}{\ohm}$, estime $\alpha$, $\omega_d$, $L$ e $C$.

\circuitoqq{%
\begin{circuitikz}[scale=0.96,transform shape,line width=0.8pt,every node/.append style={font=\scriptsize}]
\draw (0,0) to[R] (2.2,0)
      to[L] (4.5,0)
      to[C] (6.8,0) -- (6.8,-2) -- (0,-2) -- (0,0);
\node[anchor=south] at (1.1,0.25) {$R=\SI{8}{\ohm}$};
\node[anchor=south] at (3.35,0.25) {$L$};
\node[anchor=south] at (5.65,0.25) {$C$};
\draw[->] (0.9,0.70) -- (2.2,0.70);
\node[anchor=south] at (1.55,0.82) {$i(t)$};
\end{circuitikz}}{RLC em série identificado a partir de uma oscilação medida.}

\begin{figure}[H]
\centering
\begin{tikzpicture}
\begin{axis}[
 width=.88\linewidth,height=5.2cm,
 xlabel={$t$ (ms)},ylabel={$i(t)$ (u.a.)},xmin=0,xmax=32,ymin=-2.2,ymax=2.2,
 grid=both,samples=300,
 tick label style={font=\scriptsize},label style={font=\small\sffamily}]
\addplot[thick,azulprof,domain=0:32] {2*exp(-0.0433217*x)*sin(deg(0.785398*x))};
\addplot[dashed,laranja,domain=0:32] {2*exp(-0.0433217*x)};
\addplot[dashed,laranja,domain=0:32] {-2*exp(-0.0433217*x)};
\end{axis}
\end{tikzpicture}
\caption{Oscilação amortecida medida e envoltória exponencial.}
\end{figure}

\resposta{$\alpha\approx\SI{43.3}{\per\second}$;
$\omega_d\approx\SI{785.4}{\radian\per\second}$;
$L\approx\SI{92.3}{\milli\henry}$; $C\approx\SI{17.5}{\micro\farad}$}
\resolucao{A envoltória obedece $e^{-\alpha t}$. Como ela cai à metade em
$\SI{16}{\milli\second}$,
\[
\alpha=\frac{\ln2}{0{,}016}\approx\SI{43.3}{\per\second}.
\]
O período amortecido fornece
$\omega_d=2\pi/T_d\approx\SI{785.4}{\radian\per\second}$. Para RLC em série,
$\alpha=R/(2L)$, então
$L=R/(2\alpha)\approx\SI{92.3}{\milli\henry}$. Além disso,
$\omega_0=\sqrt{\omega_d^2+\alpha^2}\approx\SI{786.6}{\radian\per\second}$ e
\[
C=\frac{1}{L\omega_0^2}\approx\SI{17.5}{\micro\farad}.
\]}
\end{questao}

\gabarito[Nível 3 --- análise integrada e formas de onda]
